\chapter*{Preface}
\addcontentsline{toc}{chapter}{Preface}

Book 11's Chapter 8 previewed limits and colimits of meaning only far enough to note that failure of gluing corresponds to semantic paradox, leaving the fuller construction for the present volume. Comprehension, this book argues across its twelve chapters, requires closure: understanding is not merely the possession of many local interpretations but the existence of a universal object that makes those interpretations commute, and where no such object exists, what results is not a milder form of understanding but paradox outright. This book defines the mechanisms by which meaning aggregates across multiple scales of the RSVP plenum, using the categorical notions of limits, the intersection or consensus of meaning, and colimits, the integration or synthesis of meaning, to show how local intelligibility patches glue into global understanding, or, where they do not, how that failure itself becomes legible as a specific, measurable kind of incoherence.

\chapter*{Diagrammatic Overview}
\addcontentsline{toc}{chapter}{Diagrammatic Overview}

Every semantic structure this series has discussed, physical, cognitive, ethical, is treated in this book as a diagram in the category $\mathbf{RSVPField}$ Book 11 constructed, and its coherence is measured by whether its cone, in the limit case, or cocone, in the colimit case, exists. This book bridges category theory and hermeneutics directly: the laws of comprehension, on the account developed here, are the limits of care.

\chapter*{Reading Note}
\addcontentsline{toc}{chapter}{Reading Note}

To understand something means to identify the universal object that makes all interpretations of it commute. This is not a metaphor borrowed from category theory to dress up an independently intelligible claim about understanding; it is this book's literal thesis, and the chapters below develop the precise sense in which it is meant.

\part{Foundations of Semantic Diagrams}

